% Generated by 03_Code/make_numbers.py from 04_Results/phase5/tuning.csv. Do not edit by hand.
\begin{table*}[t]
\caption{Hyperparameter setting chosen for each arm and model from the grid of four settings per
model, by validation macro-F1 on seed 0 of the grouped split. Random forest: maximum depth and
minimum samples per leaf. LightGBM: number of leaves and learning rate. Multilayer perceptron:
hidden layer sizes and L2 penalty. UNSW-NB15 original is the release with duplicate records dropped; the
sensitivity arms of UNSW-NB15 are not listed.}\label{tab:hyper}
\small
\setlength{\tabcolsep}{4pt}
\begin{tabular}{@{}lllll@{}}
\toprule
 & & Random forest & LightGBM & Multilayer perceptron \\
Benchmark & Arm & depth / leaf & leaves / rate & layers / L2 \\
\midrule
CICIDS2017 & CIC original & none / 1 & 31 / 0.05 & (128, 64) / 0.001 \\
 & WTMC-2021 & none / 1 & 127 / 0.1 & (64) / 0.001 \\
 & CRiSIS-2022 & none / 3 & 127 / 0.05 & (128, 64) / 0.001 \\
 & LYCOS & none / 1 & 31 / 0.05 & (128, 64) / 0.001 \\
 & NFS-2023-nTE & none / 1 & 127 / 0.05 & (128, 64) / 0.0001 \\
 & NFS-2023-TE & none / 3 & 31 / 0.05 & (128, 64) / 0.001 \\
 & BCCC & none / 1 & 31 / 0.1 & (128, 64) / 0.001 \\
\midrule
UNSW-NB15 & UNSW-NB15 original & none / 1 & 31 / 0.1 & (128, 64) / 0.001 \\
 & NF-UNSW-NB15-v3 & none / 1 & 127 / 0.05 & (128, 64) / 0.0001 \\
 & CIC-UNSW-NB15 & 20 / 1 & 31 / 0.1 & (128, 64) / 0.0001 \\
\bottomrule
\end{tabular}
\end{table*}
